\chapter{Literature Review}

\section{Related Works}
The closest study to this one is a 2026 benchmark on the AgriPriceBD dataset of Bangladeshi market prices. It covers five commodities from July 2020 to June 2025 and compares naive persistence, SARIMA, Prophet, a BiLSTM, Transformer variants and Informer, all using price history only \cite{agripricebd}. No single model was best everywhere. Naive persistence was best for four of the five commodities, and Prophet did badly on step-like prices. I see the same thing in the TCB data, where prices stay still for long stretches. Some studies from other countries report that tree ensembles beat a random walk for food price indices. I have not cited them because I could not check them in enough detail, and they use index or monthly data rather than daily retail prices.

Work on onion markets in Bangladesh shows that imported onion prices strongly influence local onion prices through a non-linear, asymmetric adjustment \cite{onionnardl}. That is why the Indian onion export measures are included as a driver here, and why each driver is dated by the day it became known.

The methods used are standard ones. Automatic exponential smoothing is the ETS candidate. LightGBM is a gradient-boosted tree learner that handles mixed numeric inputs and missing values without much preparation \cite{lightgbm}, and its settings were searched with Optuna \cite{optuna}. Forecast error is reported as a ratio to a naive benchmark and in a scale-free form \cite{hk2006}. Two forecasts are compared with the Diebold--Mariano test \cite{dm1995} and its small-sample correction \cite{hln1997}. Whether the share of prices above an upper bound matches its nominal rate is checked with the proportion-of-failures test of Kupiec \cite{kupiec}. The driver breakdown uses the additive per-feature contributions of a tree model, the quantity that SHAP reports for trees \cite{shap}, summed over groups of features. Pretrained forecasting models that accept covariates, such as Chronos-2 \cite{chronos2}, exist but were not tested. Table~\ref{tab:lit} summarises these works and what each is used for here.

\begin{table}[H]\centering\small
\caption[Summary of the literature used]{Summary of the literature used in this report.}\label{tab:lit}
\begin{tabular}{p{2.6cm}p{3.6cm}p{5.2cm}}
\toprule
\textbf{Work} & \textbf{Topic} & \textbf{Relevance to this project} \\
\midrule
AgriPriceBD benchmark (2026) \cite{agripricebd} & Classical and deep models on Bangladeshi market prices, price history only & Naive persistence was hard to beat; Prophet failed on step-like prices \\
Mila et al. (2023) \cite{onionnardl} & Spatial price transmission in onion markets, NARDL & Reason for the onion export driver and for dating drivers by availability \\
Hyndman and Koehler (2006) \cite{hk2006} & Measures of forecast accuracy & Scale-free error and ratio to a naive forecast \\
Diebold and Mariano (1995) \cite{dm1995}; Harvey et al. (1997) \cite{hln1997} & Test of equal forecast accuracy and its small-sample correction & Test of the model against ``tomorrow equals today'' \\
Ke et al. (2017) \cite{lightgbm} & LightGBM & The main model \\
Akiba et al. (2019) \cite{optuna} & Optuna & Search of the model settings on 2024 \\
Kupiec (1995) \cite{kupiec} & Proportion-of-failures test & Check of the expected range \\
Lundberg and Lee (2017) \cite{shap} & SHAP, additive feature contributions & Basis of the driver breakdown \\
Ansari et al. (2025) \cite{chronos2} & Chronos-2, pretrained forecasting with covariates & Not tested; listed as future work \\
\bottomrule
\end{tabular}

\end{table}

\section{Research Gap}
Published forecasting work on Bangladeshi retail prices compares models on a few commodities using price history, and it stops at the point forecast. In the work I could find, three things were missing: a forecast with a stated range for each item of a daily retail price report; an alert rule whose level is fixed from false alarms on ordinary days and then checked against documented price episodes; and a way to say, for each alert, how much of the model's forecast comes from the measurable drivers and how much of the rise they leave unexplained. I also found no evidence that external drivers improve item-level forecasts of Bangladeshi daily retail prices.

\section{Research Contributions}
The project makes the following contributions.
\begin{itemize}
    \item A cleaned and audited daily panel of 63 TCB items over 1,412 report days, with a machine re-read of a 2\% sample against the original sheets that found no mismatch.
    \item A single-model forecasting pipeline for four horizons, compared with five alternatives on a validation year, with every setting fixed on 2024 and written to a committed file before the 2025--26 test was scored once.
    \item An expected-range and alert rule whose level was fixed from 2024 false alarms outside the event windows, checked against 12 documented episodes and against price rises defined from TCB prices alone, in both cases with a random-date control.
    \item A driver breakdown for each alert, with a stability check across refits.
    \item A clear negative result for the point forecasts: the model did not beat ``tomorrow equals today'', and the evidence is reported in full.
    \item A working set of notebooks, export files and a dashboard that run from the raw price file to the alerts and can take a longer price file later.
\end{itemize}
